\section{Automorphy inputs and the remaining specialization conditions}
\label{sec:inputs}

The congruences of Corollary~\ref{sys:congruences} connect the cohomology
of the auxiliary elliptic curve $A$ to that of the twist $B_t$ of $E$.
To turn them into a proof of automorphy, we use two results of
Caraiani--Newton. We state the precise forms needed here, and then identify
the properties that the parameter $t$ must satisfy. All the results quoted
from~\cite{CN} use the numbering of its third arXiv version.

\begin{definition}[Decomposed genericity]\label{def:decomposed-generic}
Let $L$ be a number field and let
$\overline r:G_L\to\GL_2(\overline\F_l)$ be continuous. The representation
is \emph{decomposed generic} if there is a rational prime $q\ne l$ splitting
completely in $L$ such that, for every $v\mid q$, it is unramified at $v$
and its Frobenius eigenvalues $\alpha_v,\beta_v$ satisfy
\[
 \alpha_v/\beta_v\ne q,\qquad \beta_v/\alpha_v\ne q
 \quad\text{in }\overline\F_l.
\]
\end{definition}

This is \cite[Definition~2.1.27]{CN}. Our construction will give
$q\equiv1\pmod l$, so it is enough to make the two eigenvalues distinct at
\emph{every} place above $q$.

We use potential ordinarity of Galois representations in the convention of
Section~\ref{sec:conventions}. The corresponding automorphic condition,
called \emph{$\iota$-ordinarity of weight zero}, is the ordinary Hecke
condition used in~\cite[Theorem~5.2]{CN}. It is imposed at a sufficiently deep
level and therefore allows ramified finite-order characters. We will obtain
it only for quadratic twists of Steinberg, using the local calculation in
\cite[Lemma~6.1.3(4)]{CN}.

\begin{theorem}[Caraiani--Newton, specialized form]\label{aut:lift}
Let $L$ be a CM number field and $l$ an odd prime other than $5$. Let
$r:G_L\to\GL_2(\overline\Q_l)$ be continuous, unramified outside finitely
many places, with determinant $\eps_l^{-1}$. Suppose that $r$ is potentially
semistable with labeled Hodge--Tate weights $0,1$ at every $v\mid l$,
that $\overline r$ is decomposed generic, and that
$\overline r|_{G_{L(\zeta_l)}}$ is absolutely irreducible.

Suppose there are a cuspidal weight-zero regular algebraic representation
$\pi$ of $\PGL_2(\A_L)$ and an isomorphism
$\iota:\overline\Q_l\xrightarrow{\sim}\C$ for which
$\overline{r_\iota(\pi)}\simeq\overline r$. At each $v\mid l$, assume one
of the following:
\begin{enumerate}[(i)]
\item Both local Galois representations are potentially crystalline;
one is potentially ordinary of weight zero if and only if the other is; and
$\rec_{L_v}(\pi_v)$ has zero monodromy.
\item Neither local Galois representation is potentially crystalline,
and $\pi$ is $\iota$-ordinary of weight zero at $v$.
\end{enumerate}
Then $r\simeq r_\iota(\Pi)$ for a cuspidal weight-zero regular algebraic
representation $\Pi$ of $\PGL_2(\A_L)$.
\end{theorem}

This is a sufficient-hypothesis specialization of
\cite[Theorem~5.2]{CN}. Its inverse-cyclotomic determinant hypothesis permits
$l=3$, including residual image $\SL_2(\F_3)$; see
\cite[Remark~5.2.2]{CN}. No additional residual adequacy assumption is being
made. All three applications below have $l=3$ or $l=p>100$.

\begin{theorem}[Elliptic starting point]\label{aut:seed}
Let $L$ be a CM number field not containing $\zeta_5$, and let $A/L$ be a
non-CM elliptic curve. If its mod-$5$ representation is decomposed generic
and is absolutely irreducible on $G_{L(\zeta_5)}$, then $H^1(A)$ is realized
by a cuspidal weight-zero regular algebraic representation of
$\PGL_2(\A_L)$.
\end{theorem}

This is \cite[Theorem~6.1 and Lemma~6.1.3]{CN}, after dualizing the Tate-module
convention there when necessary. Absolute irreducibility and decomposed
genericity in dimension two are preserved by duality.

For later use we also record the compatibility available after an elliptic
curve itself is known to be modular.

\begin{theorem}[Compatibility for modular elliptic curves]
\label{aut:elliptic-compat}
Let $D/L$ be a non-CM elliptic curve over a CM number field. Suppose that a
regular algebraic cuspidal representation $\pi$ realizes $H^1(D)$ at one
coefficient prime. Then $\pi$ has weight zero and trivial central character,
realizes $H^1(D)$ at every coefficient prime and every coefficient
isomorphism, and satisfies
\[
 \iota\WD\bigl(H^1(D,\overline\Q_l)|_{G_{L_v}}\bigr)^{\Fss}
 \simeq\rec^T_{L_v}(\pi_v)\qquad(v\nmid l).
\]
If $D$ has potentially multiplicative reduction at $v\mid l$, then $\pi$
is $\iota$-ordinary of weight zero there.
\end{theorem}

This is \cite[Lemma~6.1.3]{CN}. For the intermediate systems, which need not
come from elliptic curves, we use instead the general existence and
good-place compatibility of $r_\iota(\pi)$ from~\cite{HLTT,Scholze}, and
Varma's equality of the underlying semisimplified Weil representations at
$v\nmid l$~\cite[Theorem~1]{Varma}. The local reduction analysis in
Section~\ref{sec:local} will supply the additional facts needed in
Theorem~\ref{aut:lift}.

Here is what remains to be arranged for a single parameter $t\in U(K)$.
The determinant and Hodge--Tate conditions hold for every parameter by
Proposition~\ref{sys:properties}. At places above $3p$, we will make the
Jacobian of $C_n(t)$ potentially good ordinary, potentially good with all
Newton slopes $1/2$, or potentially totally toric, according to the
reduction type of $E$. This will make all four systems potentially
crystalline and potentially ordinary, potentially crystalline but not potentially ordinary, or
non-potentially crystalline, respectively. Section~\ref{sec:local} obtains
these properties from nonempty local open conditions on $t$.

Section~\ref{sec:residual} then chooses a parameter in those open sets for
which the residual hypotheses hold. The pairs $(m,l)=(3,5),(3,p),(p,3)$
provide respectively the elliptic starting point and the first two lifts;
the last lift uses the large mod-$p$ image of $E$. The two conjugates of
$K$ must be considered together to obtain decomposed genericity. Once
these properties have been proved, Section~\ref{sec:automorphy} propagates
automorphy along the three congruences and descends to $K$.
